\documentclass[a4paper]{article}
\usepackage{amsmath}
\usepackage[utf8]{inputenc}
\usepackage[czech]{babel}

\title{\vspace{-1cm}Protokol č. 1: Měření objemu tuhých těles}
\author{Michal Matiáš}
\date{\today}

\begin{document}
\maketitle

\section{Zadání}
Úkolem bylo:\\
\begin{enumerate}
\item Změřit průměr a výšku kovového válečku.
\item Na základě naměřených hodnot vypočíst kombinovanou stadnardní nejistotu jednotlivých rozměrů.
\item Vypočíst objem válečku a jeho kombinovanou standardní nejistotu.
\end{enumerate}
\section{Použité vybavení}
Pro měření průměru válečku byl použit mikrometr se stupnicí o kroku 0,01 mm, Pro měření výšky válečku bylo použito posuvné meřítko se stupnicí o kroku 0,02 mm.

\section{Naměřené hodnoty}
V tabulce \ref{table:1} jsou uvedeny naměřené hodnoty pro průměr d a výšku h.
\begin{center}
\begin{table}[h!]
\begin{tabular}{|| l | l | l | l | l | l | l | l | l | l | l ||}
\hline
Meření č. & 1 & 2 & 3 & 4 & 5 & 6 & 7 & 8 & 9 & 10\\
\hline\hline
d [mm] & 14,79 & 14,84 & 14,82 & 14,81 & 14,80 & 14,90 & 14,93 & 14,94 & 14,82 & 14,85\\
\hline
h [mm] & 37,40 & 37,36 & 37,44 & 37,38 & 37,44 & 37,40 & 37,38 & 37,40 & 37,36 & 37,42\\
\hline
\end{tabular}
\caption{Nameřené hodnoty pro průměr d a výšku h válečku}
\label{table:1}
\end{table}
\end{center}

\section{Výpočet}
Výpočet objemu a nejistot byl vypočten pomocí následujícího matlabovského skriptu, ve kterém \verb|uAx|, \verb|uBx| a \verb|uCx| jsou standardní nejistoty typu A, B a kombinovaná nejistota příslušné veličiny a \verb|V_mean| je výsledný objem:
\newpage
\begin{verbatim}
clear; home;
d = [14.79 14.84 14.82 14.81 14.80 14.90 14.93 14.94 14.82 14.85]
h = [37.40 37.36 37.44 37.38 37.44 37.40 37.38 37.40 37.36 37.42]

uAd = sqrt(1/(length(d)*(length(d) - 1))*sum((d - mean(d)) .^ 2))
uAh = sqrt(1/(length(h)*(length(h) - 1))*sum((h - mean(h)) .^ 2))

uBd = 0.01/sqrt(12)
uBh = 0.02/sqrt(12)

uCd = sqrt(uAd^2 + uBd^2)
uCh = sqrt(uAh^2 + uBh^2)

uCV = sqrt((1/2*pi*mean(d)*mean(h)*uCd)^2+(1/4*mean(d)^2*pi*uCh)^2)

V_mean = 1/4*pi*mean(d)^2*mean(h)
\end{verbatim}

\section{Závěr}
Naměřené hodnoty jsou uvedeny v tabulce \ref{table:1}. Po odečtení hodnot z matlabu a stanovení výsledku na dvě platné cifry dostáváme:\\
$u_{C_d} = 0,017 mm$\\
$u_{C_h} = 0,011 mm$\\
$u_{C_V} = 15 mm$\\
Můžeme psát:\\
$\mathbf{d = (14,85 \pm 0,017)mm}$\\
$\mathbf{h = (37,398 \pm 0,011)mm}$\\
$\mathbf{V = (6477 \pm 15)mm^3}$

\end{document}
